\chapter{Thứ còn lại}
\markboth{Thứ còn lại}{What Remains}

\section{Một người qua biến cố mới biết điều còn lại là ai.}
\begin{truyen}{Một người qua biến cố mới biết điều còn lại là ai.}{What remains reveals what mattered.}
\chuhoa{S}{au một lần ngã mạnh, anh mất tiền, mất kế hoạch, mất nhiều mối quen. Nhưng vài người vẫn ở lại. Và lần đầu tiên anh thấy rõ thứ thật sự gọi là tài sản của đời mình.}
\textit{(After a hard fall, he lost money, plans, and many loose connections. But a few people stayed. For the first time, he clearly saw what truly counted as wealth in his life.)}
\end{truyen}
\begin{baihoc}
Mất đi nhiều lớp ngoài cùng đôi khi mới lộ ra điều cốt lõi mình thật sự có.
\textit{(Only when many outer layers are lost does the core of what we truly have appear.)}
\end{baihoc}
\begin{ghinhoanh}\textbf{Short English to remember:}

What remains shows what mattered.
\end{ghinhoanh}
\begin{tuhocnhanh}
1. Nếu mất nhiều thứ ngoài kia, điều gì trong đời em vẫn còn? \textit{(If you lost many outer things, what would still remain in your life?)}

2. Em có đang nuôi đủ điều còn lại ấy không? \textit{(Are you caring enough for those remaining essentials?)}
\end{tuhocnhanh}
\ngancach

\section{Một người hiểu rằng mất mặt không đáng sợ bằng mất nhân cách.}
\begin{truyen}{Một người hiểu rằng mất mặt không đáng sợ bằng mất nhân cách.}{Face can return; character may not.}
\chuhoa{C}{ó lúc anh xấu hổ vì thừa nhận sai. Nhưng rồi anh thấy điều đáng sợ hơn là giữ sĩ diện bằng cách tiếp tục sai.}
\textit{(At times he felt ashamed to admit he was wrong. But he came to see that what was more frightening was protecting face by continuing in the wrong.)}
\end{truyen}
\begin{baihoc}
Mất chút thể diện để giữ lấy nhân cách là một cái mất rất nên chấp nhận.
\textit{(Losing a little face to preserve character is a loss well worth accepting.)}
\end{baihoc}
\begin{ghinhoanh}\textbf{Short English to remember:}

Face can return; character may not.
\end{ghinhoanh}
\begin{tuhocnhanh}
1. Em sợ mất mặt hay sợ sống sai hơn? \textit{(Do you fear losing face more than living wrongly?)}

2. Em cần dũng cảm sửa điều gì? \textit{(What do you need courage to correct?)}
\end{tuhocnhanh}
\ngancach

\section{Một người mất hết kế hoạch nhưng được lại chính mình.}
\begin{truyen}{Một người mất hết kế hoạch nhưng được lại chính mình.}{Broken plans can uncover a truer life.}
\chuhoa{K}{hi mọi dự định đổ xuống, cô tưởng mình trắng tay. Nhưng cũng chính lúc ấy, cô ngừng chạy theo khuôn người khác đặt, và bắt đầu sống giống mình hơn.}
\textit{(When every plan collapsed, she thought she had nothing left. Yet it was exactly then that she stopped chasing others' expectations and began to live more like herself.)}
\end{truyen}
\begin{baihoc}
Không phải đổ vỡ nào cũng cướp đi; có cái đổ vỡ gỡ ta ra khỏi một đời không thuộc về mình.
\textit{(Not every collapse only takes away; some free us from a life that was never truly ours.)}
\end{baihoc}
\begin{ghinhoanh}\textbf{Short English to remember:}

Broken plans can reveal a truer life.
\end{ghinhoanh}
\begin{tuhocnhanh}
1. Có kế hoạch nào đổ vỡ nhưng lại dạy em điều đúng hơn? \textit{(Has any broken plan taught you something truer?)}

2. Em có đang sống đời của mình hay của kỳ vọng khác? \textit{(Are you living your own life or someone else's expectation?)}
\end{tuhocnhanh}
\ngancach

\section{Một người hiểu rằng cái còn lại sau cùng là cách mình đã sống.}
\begin{truyen}{Một người hiểu rằng cái còn lại sau cùng là cách mình đã sống.}{In the end, your manner of living remains.}
\chuhoa{C}{àng lớn, anh càng thấy của cải, lời khen, vị trí đều đổi. Nhưng cách mình đối xử với người khác, cách mình giữ lời, cách mình đi qua những lúc tối mới là thứ đọng lại lâu.}
\textit{(The older he got, the more he saw wealth, praise, and status all changing. But how he treated others, kept his word, and walked through dark times were what truly remained.)}
\end{truyen}
\begin{baihoc}
Thứ ở lại sâu nhất không phải ta đã có gì, mà là ta đã là người thế nào.
\textit{(What remains most deeply is not what we had, but what kind of person we were.)}
\end{baihoc}
\begin{ghinhoanh}\textbf{Short English to remember:}

How you live remains longest.
\end{ghinhoanh}
\begin{tuhocnhanh}
1. Nếu nhìn lại, em muốn đời mình được nhớ bởi điều gì? \textit{(If looking back, how would you want your life to be remembered?)}

2. Hôm nay em đã sống gần điều đó chưa? \textit{(Have you lived near that today?)}
\end{tuhocnhanh}
\ngancach

\section{Một người hiểu muộn rằng mất để lớn cũng là một dạng được.}
\begin{truyen}{Một người hiểu muộn rằng mất để lớn cũng là một dạng được.}{Some losses arrive as disguised gifts.}
\chuhoa{S}{au một lần hụt điều mình rất muốn, cô khổ sở rất lâu. Vài năm sau, cô mới hiểu nếu khi ấy mình không mất, có lẽ mình cũng không lớn nổi theo cách hôm nay.}
\textit{(After losing something she deeply wanted, she suffered for a long time. Years later, she understood that if she had not lost it then, she might never have grown into who she is now.)}
\end{truyen}
\begin{baihoc}
Mất mát không phải lúc nào cũng là kẻ cướp; đôi khi nó là thầy giáo khó tính của đời.
\textit{(Loss is not always a thief; sometimes it is life's strict teacher.)}
\end{baihoc}
\begin{ghinhoanh}\textbf{Short English to remember:}

Some losses are disguised gifts.
\end{ghinhoanh}
\begin{tuhocnhanh}
1. Mất mát nào của em về sau lại hóa ra có ích? \textit{(What loss in your life later turned out to be useful?)}

2. Em học được gì từ điều đó? \textit{(What did you learn from it?)}
\end{tuhocnhanh}
\ngancach